
\documentclass{article}
\usepackage{colortbl}
\usepackage{makecell}
\usepackage{multirow}
\usepackage{supertabular}

\begin{document}

\newcounter{utterance}

\centering \large Interaction Transcript for game `hot\_air\_balloon', experiment `air\_balloon\_survival\_en\_negotiation\_hard', episode 0 with Qwen3{-}32B{-}t0.0.
\vspace{24pt}

{ \footnotesize  \setcounter{utterance}{1}
\setlength{\tabcolsep}{0pt}
\begin{supertabular}{c@{$\;$}|p{.15\linewidth}@{}p{.15\linewidth}p{.15\linewidth}p{.15\linewidth}p{.15\linewidth}p{.15\linewidth}}
    \# & $\;$A & \multicolumn{4}{c}{Game Master} & $\;\:$B\\
    \hline

    \theutterance \stepcounter{utterance}  
    & & \multicolumn{4}{p{0.6\linewidth}}{
        \cellcolor[rgb]{0.9,0.9,0.9}{
            \makecell[{{p{\linewidth}}}]{
                \texttt{\tiny{[P1$\langle$GM]}}
                \texttt{You are participating in a collaborative negotiation game.} \\
\\ 
\texttt{Together with another participant, you must agree on a single set of items that will be kept. Each of you has your own view of how much each item matters to you (importance). You do not know how the other participant values the items. Additionally, you are given the effort each item demands.} \\
\texttt{You may only agree on a set if the total effort of the selected items does not exceed a shared limit:} \\
\\ 
\texttt{LIMIT = 3624} \\
\\ 
\texttt{Here are the individual item effort values:} \\
\\ 
\texttt{Item effort = \{"C17": 210, "B41": 291, "A71": 831, "C00": 111, "B35": 926, "A77": 827, "A03": 25, "A04": 121, "C54": 583, "C01": 766, "A87": 14, "B32": 559, "B86": 304, "B33": 989, "A10": 691\}} \\
\\ 
\texttt{Here is your personal view on the importance of each item:} \\
\\ 
\texttt{Item importance values = \{"C17": 138, "B41": 583, "A71": 868, "C00": 822, "B35": 783, "A77": 65, "A03": 262, "A04": 121, "C54": 508, "C01": 780, "A87": 461, "B32": 484, "B86": 668, "B33": 389, "A10": 808\}} \\
\\ 
\texttt{Goal:} \\
\\ 
\texttt{Your goal is to negotiate a shared set of items that benefits you as much as possible (i.e., maximizes total importance to YOU), while staying within the LIMIT. You are not required to make a PROPOSAL in every message {-} you can simply negotiate as well. All tactics are allowed!} \\
\\ 
\texttt{Interaction Protocol:} \\
\\ 
\texttt{You may only use the following structured formats in a message:} \\
\\ 
\texttt{PROPOSAL: \{'A', 'B', 'C', …\}} \\
\texttt{Propose keeping exactly those items.} \\
\texttt{REFUSE: \{'A', 'B', 'C', …\}} \\
\texttt{Explicitly reject opponent's proposal.} \\
\texttt{ARGUMENT: \{'...'\}} \\
\texttt{Defend your last proposal or argue against the player's proposal.} \\
\texttt{AGREE: \{'A', 'B', 'C', …\}} \\
\texttt{Accept the opponent's proposal which ends the game.} \\
\texttt{STRATEGIC REASONING: \{'...'\}} \\
\texttt{	Describe your strategic reasoning or anticipation explaining your choice of action. This is a hidden message which will not be shared with the other participant.} \\
\\ 
\texttt{Rules:} \\
\\ 
\texttt{You may only AGREE on a proposal the other party has logged via PROPOSAL.} \\
\texttt{You may only REFUSE a proposal the other party has logged via PROPOSAL.} \\
\texttt{Total effort of any PROPOSAL or AGREE set must be ≤ LIMIT.} \\
\texttt{Do NOT reveal your hidden importance scores.} \\
\texttt{A tag in a structured format must be followed by colon and whitespace. The argument must be a python set containing 0 or more strings.} \\
\texttt{So, it must be of the form TAG: \{...\}} \\
\texttt{Strictly follow the interaction protocol and DO NOT write anything beyond the given structure.} \\
\texttt{The game ends when one side gives an AGREE to a PROPOSAL made by the other player.} \\
\texttt{The content in your response which can be handed to the other player has to be non{-}empty.} \\
\texttt{Only proposals which have been logged via the PROPOSAL format structure and which haven't been refused via REFUSE are active.} \\
\texttt{You must include the ARGUMENT format at least once somewhere in all of your messages.} \\
\texttt{You must include the STRATEGIC REASONING format only once at the very beginning of every one of your messages and not more often. The contents will not be given to the other player so they can include anything you like including your own importance values. Here you should reason step by step to come up with you next move.} \\
            }
        }
    }
    & \\ \\

    \theutterance \stepcounter{utterance}  
    & \multicolumn{4}{p{0.6\linewidth}}{
        \cellcolor[rgb]{0.9,0.9,0.9}{
            \makecell[{{p{\linewidth}}}]{
                \texttt{\tiny{[P1$\rangle$GM]}}
                \texttt{<think>} \\
\texttt{Okay, let's see. I need to figure out the best set of items to propose that maximizes my total importance without exceeding the effort limit of 3624. First, I should calculate the total effort for different combinations. But since I don't know the other person's preferences, I have to focus on my own importance values.} \\
\\ 
\texttt{Looking at the items, I should prioritize those with high importance and low effort. Let me list out the items with their importance and effort:} \\
\\ 
\texttt{{-} A87: Importance 461, Effort 14 (super low effort, high value)} \\
\texttt{{-} C00: 822 importance, 111 effort} \\
\texttt{{-} A71: 868, 831} \\
\texttt{{-} A10: 808, 691} \\
\texttt{{-} B41: 583, 291} \\
\texttt{{-} B86: 668, 304} \\
\texttt{{-} C54: 508, 583} \\
\texttt{{-} B32: 484, 559} \\
\texttt{{-} C17: 138, 210} \\
\texttt{{-} B35: 783, 926 (high importance but high effort)} \\
\texttt{{-} A77: 65, 827 (low importance, high effort)} \\
\texttt{{-} A03: 262, 25} \\
\texttt{{-} A04: 121, 121} \\
\texttt{{-} B33: 389, 989 (high effort)} \\
\texttt{{-} C01: 780, 766} \\
\\ 
\texttt{Hmm. A87 is a no{-}brainer because it's super low effort and high importance. Then A03 is also low effort. Let's start by including A87 and A03. Their combined effort is 14 +25 = 39. That leaves 3624 {-}39 = 3585.} \\
\\ 
\texttt{Next, the highest importance items are A71 (868) and A10 (808). Let's check their efforts. A71 is 831, A10 is 691. If I take both, that's 831 +691 = 1522. Adding to the previous total:} \\
            }
        }
    }
    & & \\ \\

    \theutterance \stepcounter{utterance}  
    & & \multicolumn{4}{p{0.6\linewidth}}{
        \cellcolor[rgb]{0.9,0.9,0.9}{
            \makecell[{{p{\linewidth}}}]{
                \texttt{\tiny{[P1$\langle$GM]}}
                \texttt{Your response did not start with the proper strategic reasoning tag at the very beginning of your response. The very first structured format must be of the form STRATEGIC REASONING: \{...\}. Try again.} \\
            }
        }
    }
    & \\ \\

    \theutterance \stepcounter{utterance}  
    & \multicolumn{4}{p{0.6\linewidth}}{
        \cellcolor[rgb]{0.9,0.9,0.9}{
            \makecell[{{p{\linewidth}}}]{
                \texttt{\tiny{[P1$\rangle$GM]}}
                \texttt{<think>} \\
\texttt{Okay, I need to make sure I start with the STRATEGIC REASONING tag. Let me think through this again. My goal is to maximize my total importance without exceeding the effort limit. Let me list out the items with their importance and effort to prioritize.} \\
\\ 
\texttt{First, A87 has the highest importance per effort ratio (461/14) and is very low effort. Then A03 is also low effort with decent importance. Including these two gives me a good start. Next, the highest importance items are A71 (868) and A10 (808). Their efforts are 831 and 691. If I take both, that's 1522 effort. Adding A87 and A03's 39, total is 1561. That leaves 3624 {-}1561 = 2063. } \\
\\ 
\texttt{Next, B41 (583, 291) and B86 (668, 304) are next in importance. Adding B41 would add 291 effort, total becomes 1561 +291 = 1852. Then B86 adds 304, total 2156. Still under the limit. Then C00 (822, 111) is high importance. Adding that would be 2156 +111 = 2267. Then C54 (508, 583) adds 583, total 2850. Still under. Then B32 (484, 559) adds 559, total 3409. Then C17 (138, 210) adds 210, total 3619. That's under 3624. } \\
\\ 
\texttt{Wait, let me check the total effort: 14 (A87) +25 (A03) +831 (A71) +691 (A10) +291 (B41) +304 (B86) +111 (C00) +583 (C54) +559 (B32} \\
            }
        }
    }
    & & \\ \\

    \theutterance \stepcounter{utterance}  
    & & \multicolumn{4}{p{0.6\linewidth}}{
        \cellcolor[rgb]{0.9,0.9,0.9}{
            \makecell[{{p{\linewidth}}}]{
                \texttt{\tiny{[P1$\langle$GM]}}
                \texttt{Your response did not start with the proper strategic reasoning tag at the very beginning of your response. The very first structured format must be of the form STRATEGIC REASONING: \{...\}. Try again.} \\
            }
        }
    }
    & \\ \\

    \theutterance \stepcounter{utterance}  
    & \multicolumn{4}{p{0.6\linewidth}}{
        \cellcolor[rgb]{0.9,0.9,0.9}{
            \makecell[{{p{\linewidth}}}]{
                \texttt{\tiny{[P1$\rangle$GM]}}
                \texttt{<think>} \\
\texttt{Okay, I need to make sure I start with the STRATEGIC REASONING tag. Let me outline my plan step by step.} \\
\\ 
\texttt{First, I'll calculate the total effort for the items I value most. A87 has the highest importance per effort (461/14) and is very low effort. A03 is also low effort with decent importance. Including these two gives a good start. Next, the highest importance items are A71 (868) and A10 (808). Their efforts are 831 and 691. Taking both adds 1522 effort. Adding A87 and A03's 39, total is 1561. Remaining effort: 3624 {-}1561 = 2063.} \\
\\ 
\texttt{Next, B41 (583, 291) and B86 (668, 304) are next in importance. Adding B41 adds 291, total 1852. Then B86 adds 304, total 2156. Still under. C00 (822, 111) adds 111, total 2267. C54 (508, 583) adds 583, total 2850. B32 (484, 559) adds 559, total 3409. C17 (138, 210) adds 210, total 3619. That's under the limit. } \\
\\ 
\texttt{Wait, let me check the total effort again: 14 +25 +831 +691 +291 +304 +111 +583 +559 +210 = 3619. Yes, under 3624. This set includes A87, A03, A71, A10, B41, B86, C00, C54, B32, C17. Total importance: 461 +262 +868 +808 +583 +668 +822 +508 +} \\
            }
        }
    }
    & & \\ \\

    \theutterance \stepcounter{utterance}  
    & & & \multicolumn{2}{p{0.3\linewidth}}{
        \cellcolor[rgb]{0.9,0.9,0.9}{
            \makecell[{{p{\linewidth}}}]{
                \texttt{\tiny{[GM$|$GM]}}
                \texttt{abort} \\
            }
        }
    }
    & & \\ \\

    \theutterance \stepcounter{utterance}  
    & & & \multicolumn{2}{p{0.3\linewidth}}{
        \cellcolor[rgb]{0.9,0.9,0.9}{
            \makecell[{{p{\linewidth}}}]{
                \texttt{\tiny{[GM$|$GM]}}
                \texttt{end game} \\
            }
        }
    }
    & & \\ \\

\end{supertabular}
}

\end{document}
